\subsection{Безразмерная постановка}

Переход к безразмерным переменным унифицирует три частные постановки, сформулированные в подразделах~2.1.2-2.1.4, и позволяет сравнивать различные конфигурации подшипников и режимы их работы на единой шкале. Кроме того, безразмеризация сокращает число определяющих параметров задачи до нескольких безразмерных групп, что упрощает параметрический анализ. Введённый аппарат используется далее при исследовании влияния микроструктур поверхности (раздел~2.2) и при численной дискретизации (глава~3).

\subsubsection*{Характерные масштабы и безразмерные переменные}

Общая схема масштабирования повторяет подход, применённый при выводе уравнения Рейнольдса в подразделе~2.1.1~\cite{hori2002,szeri1998}, но теперь записывается в обобщённых обозначениях, пригодных для всех трёх постановок:
%
\begin{equation}
\bar{s}_1 = \frac{s_1}{L_1}, \quad
\bar{s}_2 = \frac{s_2}{L_2}, \quad
\bar{h}   = \frac{h}{h_0},   \quad
\bar{p}   = \frac{p}{p_0},   \quad
\bar{U}   = \frac{U}{U_0},   \quad
\bar{t}   = \frac{t}{t_0},
\label{eq:nondim_vars}
\end{equation}
%
\begin{conditions}
$s_1$, $s_2$ -- поверхностные координаты (конкретные для каждой постановки, см.\ таблицу~\ref{tab:nondim_scales} ниже);\\
$h_0$ -- характерная толщина зазора (радиальный зазор $c$ или средний зазор);\\
$U_0$ -- характерная скорость увлечения;\\
$t_0$ -- характерное время (для нестационарных задач).
\end{conditions}
Масштаб давления выбирается так, чтобы клиновой член в правой части уравнения Рейнольдса имел порядок единицы. Из баланса~\eqref{eq:pressure_scale} следует:
%
\begin{equation}
p_0 = \frac{6\,\mu_0\,U_0\,L_1}{h_0^2},
\label{eq:nondim_p0}
\end{equation}
%
\begin{conditions}
$L_1$ -- масштаб длины вдоль направления увлечения (координата $s_1$);\\
$\mu_0$ -- характерное значение вязкости.
\end{conditions}

\subsubsection*{Безразмерное уравнение Рейнольдса}

Подставляя масштабы~\eqref{eq:nondim_vars},~\eqref{eq:nondim_p0} в обобщённое уравнение Рейнольдса~\eqref{eq:reynolds_full} и сокращая размерные множители, получаем единую безразмерную форму:
%
\begin{equation}
\frac{\partial}{\partial \bar{s}_1}\!\left(\bar{h}^{\,3}\,\frac{\partial \bar{p}}{\partial \bar{s}_1}\right)
+ \Lambda^2\,\frac{\partial}{\partial \bar{s}_2}\!\left(\bar{h}^{\,3}\,\frac{\partial \bar{p}}{\partial \bar{s}_2}\right)
= \bar{U}\,\frac{\partial \bar{h}}{\partial \bar{s}_1}
+ \sigma\,\frac{\partial \bar{h}}{\partial \bar{t}},
\label{eq:reynolds_nondim}
\end{equation}
%
где введены две безразмерные группы:
%
\begin{equation}
\Lambda = \frac{L_1}{L_2}
\label{eq:nondim_Lambda}
\end{equation}
%
-- параметр удлинённости (aspect ratio), характеризующий соотношение продольного и поперечного масштабов, и
%
\begin{equation}
\sigma = \frac{2\,L_1}{U_0\,t_0}
\label{eq:nondim_sigma}
\end{equation}
%
-- параметр нестационарности (squeeze number), определяющий относительный вклад выжимания плёнки.

При $\sigma = 0$ (либо в стационарной задаче) член выжимания отсутствует, и уравнение содержит только клиновой источник. Параметр $\Lambda$ определяет относительный вклад поперечного растекания: при $\Lambda \ll 1$ вклад второго члена мал и задача приближается к одномерной по $s_1$; при $\Lambda \gg 1$ доминирует поперечное растекание по $s_2$, что соответствует <<короткому>> приближению в терминологии радиального подшипника (подраздел~2.1.3). Для цилиндрических координат, используемых в постановках~2.1.2-2.1.4, в левой части уравнения~\eqref{eq:reynolds_nondim} появляются дополнительные множители $1/\bar{r}$ и $1/\bar{r}^{\,2}$, однако общая структура уравнения сохраняется.

\subsubsection*{Масштабы для частных постановок}

Конкретные значения масштабов для трёх постановок, рассмотренных в подразделах~2.1.2-2.1.4, приведены в таблице~\ref{tab:nondim_scales}.

\begin{table}[!ht]
\centering
\caption{Характерные масштабы для безразмерной постановки}
\label{tab:nondim_scales}
\resizebox{\textwidth}{!}{%
\begin{tabular}{|l|c|c|c|c|c|c|c|c|}
\hline
\multicolumn{1}{|c|}{\textbf{Постановка}} & $\boldsymbol{s_1}$ & $\boldsymbol{s_2}$ & $\boldsymbol{L_1}$ & $\boldsymbol{L_2}$ & $\boldsymbol{h_0}$ & $\boldsymbol{U_0}$ & $\boldsymbol{p_0}$ & $\boldsymbol{t_0}$ \\
\hline
Упорная (2.1.2)
  & $R\theta$ & $r$ & $R$ & $R_{\mathrm{out}} - R_{\mathrm{in}}$
  & $h_{\min}$ или $h_{\mathrm{mean}}$
  & $\omega R$ & $\dfrac{6\mu_0\omega R^2}{h_0^2}$ & --
\\[6pt]
\hline
Радиальная (2.1.3)
  & $R\theta$ & $z$ & $R$ & $L$
  & $c$
  & $\omega R$ & $\dfrac{6\mu_0\omega R^2}{c^2}$ & --
\\[6pt]
\hline
Возвратно-пост.\ (2.1.4)
  & $z$ & $R\theta$ & $L$ & $R$
  & $c$
  & $\max|U(t)|$ & $\dfrac{6\mu_0\,U_0\,L}{c^2}$ & $1/\Omega$
\\[6pt]
\hline
\end{tabular}%
}
\end{table}

\FloatBarrier

В упорной и радиальной постановках задача стационарна, поэтому временной масштаб $t_0$ не требуется (обозначено <<-->>{} в таблице). Координата $s_1$ соответствует направлению увлечения, а $s_2$ -- поперечному направлению. Для упорного подшипника $R$ -- характерный радиус (например, $R_{\mathrm{out}}$ или среднее значение $(R_{\mathrm{in}} + R_{\mathrm{out}})/2$). При подстановке масштабов из таблицы в общее уравнение~\eqref{eq:reynolds_nondim} получаются безразмерные формы, согласованные с размерными уравнениями~\eqref{eq:thrust_reynolds_num},~\eqref{eq:journal_reynolds_num},~\eqref{eq:recip_reynolds_num}.

Характерный масштаб давления~$p_0 = 6\mu_0 U_0 L_1 / h_0^2$
(формула~\eqref{eq:nondim_p0}) определяет, при каких реальных
давлениях безразмерное поле~$P$ имеет порядок единицы.
Для радиального подшипника насоса (глава~5): $\mu_0 = 0{,}01105$~Па$\cdot$с,
$U_0 = \omega R = 10{,}9$~м/с, $L_1 = R = 0{,}035$~м,
$h_0 = c = 50$~мкм:
%
\begin{equation}
  p_0 = \frac{6 \cdot 0{,}01105 \cdot 10{,}9
    \cdot 0{,}035}{(50 \cdot 10^{-6})^2}
  \approx 10{,}1 \;\text{МПа}.
  \label{eq:p0_example}
\end{equation}
%
Максимальные безразмерные давления в расчётах глав 5--6 имеют
порядок $10^{-1}$, что соответствует размерным значениям
порядка единиц мегапаскалей,~-- в пределах применимости
изотермической постановки. Для упорного подшипника (глава~7)
масштаб~$p_0$ иной вследствие меньшего характерного зазора
и другой скорости, что объясняет различие в абсолютных
значениях несущей способности при схожих безразмерных уровнях
давления.

Параметр $\Lambda = L_1/L_2$~\eqref{eq:nondim_Lambda}
физически характеризует соотношение «продольного»
и «поперечного» масштабов задачи вдоль направления увлечения
и поперёк него (таблица~\ref{tab:nondim_scales}).
Для радиального подшипника $L_1 = R$, $L_2 = L$ (длина
подшипника), поэтому $\Lambda = R/L$. При малых $\Lambda$
($R \ll L$, узкий и длинный подшипник) поперечным растеканием
можно пренебречь, и задача сводится к одномерной по~$s_1$~--
это «бесконечно длинное» приближение. При больших $\Lambda$
($R \gg L$, широкий и короткий подшипник) давление быстро
выравнивается по~$s_1$, и поперечный градиент определяет
структуру течения~-- это «короткое» приближение. Для реальных
подшипников $L/D \approx 0{,}5$--$1{,}5$, откуда
$\Lambda \approx 0{,}33$--$1{,}0$, что соответствует
промежуточному режиму с существенным вкладом обоих членов
в уравнении~\eqref{eq:reynolds_nondim}.

Введённые масштабы определяют безразмерную форму уравнения Рейнольдса и базовой геометрии зазора. Параметры, характеризующие микроструктуры поверхности, -- безразмерная глубина $h_p/h_0$, размеры элементов, плотность покрытия -- вводятся отдельно в подразделе~2.2.4. Далее в работе, если не оговорено иное, используются безразмерные переменные; черта над символом опускается для краткости записи.
